\section{Implementation and measurements}\label{sec:empirical}

We measure the C/GMP implementation \texttt{intmul\_fastdft} (\texttt{algo=1}, about $1{,}700$ lines of C~\cite{CGMP26}) of the algorithm of this paper, with the windowed sinc-Gaussian kernel and the additive schedule, and with inner integer products delegated to GMP (it does not implement the companion paper's $d$-dimensional algorithm~\cite{CLW26}). Compared with the Turing machine specification in Sections~\ref{sec:prelim}--\ref{sec:main-proof}, \texttt{intmul\_fastdft.c} uses the identical additive schedule~\eqref{eq:cert-sched} ($p_{\mathrm{work},\ell} = p_\ell + K_\ell + 6$, $p_{\ell+1} = p_\ell + 2 K_\ell + 12$, $p_0 = 2b + 2K_0 + 6$), evaluates precomputed tables with MPFR nearest rounding, packs signed Kronecker slots into non-negative GMP \texttt{mpz\_mul} operands via a precomputed bias offset, performs direct fixed-point division in the two-term column solver, and lowers the default leaf cutoff to $\texttt{bt} = p_0^2$ evaluated at the root $p_0$ (rather than $4 p_\ell^2$). At $n = 10^9$ bits ($b = 29$, $K_0 = 27$, $R_0 = 2^{27}$, $p_0 = 118$), $\texttt{bt} = 13{,}924 < R_1 = 2^{14} = 16{,}384$, so the $10^9$-bit run exercises two recursion levels ($2^{27} \to 2^{14} \to 2^7 = 128$). At that level $R_1 = 2^{14} < 4 p_1^2 = 135{,}424$, so these nodes lie outside the non-leaf hypothesis of Theorem~\ref{thm:spread} and Lemma~\ref{lem:forward}(1) (and exercise the folding case $W_1 = 191 > r_1 = 128$ of Remark~\ref{rem:folding}); the bit-exact agreement observed at $10^9$ bits is therefore an empirical observation, not a consequence of the proofs. The runs at $10^7$, $10^8$, and $5\cdot 10^8$ bits treat level~$1$ as a leaf and stay within the parameter regime of the proofs (the implementation differences listed above remain).

All measurements were run on an AMD EPYC 7B13 (Zen~3) virtual machine under KVM ($64$ vCPUs, $117\,\mathrm{GiB}$ RAM, Linux x86\_64), compiled with \texttt{gcc}~15.2.0 (\texttt{-O2 -march=native -fopenmp -DNDEBUG}) against GMP~6.3.0 and MPFR~4.2.2~\cite{CGMP26}. The measured binary was compiled from a copy of \texttt{intmul\_fastdft.c} that differs from the shipped file only in a header comment and two preprocessor-guard blocks (an unused include branch under a build-internal macro, removed from the shipped file, and a guard added so that the shipped file also compiles without OpenMP; under the stated flags neither changes the compiled code), with both SHA-256 digests recorded in \texttt{bench/MACHINE.txt}; the host was lightly shared (load average $2$--$5$ on $64$ vCPUs) during the single-threaded runs, and `one core' means one vCPU of the KVM guest. Each run (one random operand pair per size) invoked \texttt{intmul\_fastdft mul <n> 0 7 <algo>} (\texttt{seed}~7, random $n$-bit operands with top bit set) and compared the result bit for bit against single-threaded GMP \texttt{mpz\_mul} on the same operands in the same process. The program's built-in forced-deep self-test at $4{,}096$ and $65{,}536$ bits reported \texttt{SELFTEST: ALL PASSED} (every product \texttt{EXACT}, both kernels), and every benchmark run from $10^7$ to $10^9$ bits printed \texttt{EXACT MATCH} (bit-exact on every tested input).

\begin{table}[!htbp]
\centering
\footnotesize
\caption{Single-threaded measurements (\texttt{OMP\_NUM\_THREADS=1}) of the C/GMP implementation \texttt{intmul\_fastdft} (\texttt{algo=1}, windowed sinc-Gaussian kernel and additive schedule; inner products by GMP) against single-threaded GMP \texttt{mpz\_mul} on the same operands (\texttt{seed}~7, AMD EPYC 7B13). Every run printed \texttt{EXACT MATCH}. Slowdown ratios are those printed by the program from the unrounded timings; recomputing them from the rounded timings shown here can differ by one in the last digit.}
\label{tab:c-benchmarks}
\setlength{\tabcolsep}{3.0pt}
\begin{tabular}{@{}llrrrrrr@{}}
\toprule
$n$ (bits) & Chain $(R_\ell,\,p_\ell)$ & Core (s) & End-to-end (s) & Tables (s) & \texttt{mpz\_mul} (s) & Slowdown & Peak RSS \\
\midrule
$10^7$ & $2^{20}\,(92) \to 1024\,(144)$ & $9.63$ & $9.92$ & $1.82$ & $0.0478$ & $207\times$ & $0.56\,\mathrm{GiB}$ \\
$10^8$ & $2^{23}\,(104) \to 4096\,(162)$ & $90.00$ & $92.51$ & $5.46$ & $0.6059$ & $153\times$ & $4.4\,\mathrm{GiB}$ \\
$5\cdot 10^8$ & $2^{26}\,(114) \to 8192\,(178)$ & $786.57$ & $805.37$ & $19.44$ & $3.4417$ & $234\times$ & $33.8\,\mathrm{GiB}$ \\
$10^9$ & $2^{27}\,(118) \to 2^{14}\,(184) \to 128\,(224)$ & $2574.67$ & $2612.65$ & $25.95$ & $8.0969$ & $323\times$ & $67.3\,\mathrm{GiB}$ \\
\bottomrule
\end{tabular}
\end{table}

Table~\ref{tab:c-benchmarks} gives the single-threaded timings (\texttt{OMP\_NUM\_THREADS=1}). Here ``core'' is the time in the two top-level \texttt{FastDFT} calls, ``end-to-end'' includes chunk unpacking, both transforms, pointwise squaring, and carry assembly, and ``tables'' is the one-time filter and twiddle setup for that size (including table setup, the single-threaded slowdowns relative to \texttt{mpz\_mul} are $245\times$, $162\times$, $240\times$, and $326\times$ as reported in \texttt{bench/\allowbreak BENCHMARK.md}; recomputed from the rounded timings of Table~\ref{tab:c-benchmarks} the first is $246\times$). Single-threaded on one Zen~3 core, the C/GMP program multiplies two $10^8$-bit integers in $92.5\,\mathrm{s}$, about $153$ times slower than GMP's \texttt{mpz\_mul} ($0.606\,\mathrm{s}$) on the same operands; it agrees bit for bit with GMP at every size tested, up to $10^9$ bits. The $10^9$-bit run needs about $72\,\mathrm{GB}$ of RAM (peak RSS $67.3\,\mathrm{GiB}$). In one $64$-thread run at $10^8$ bits (\texttt{OMP\_NUM\_THREADS=64}, same binary and operands), \texttt{intmul\_fastdft} (\texttt{algo=1}) took $4.90\,\mathrm{s}$ in \texttt{FastDFT} core and $6.31\,\mathrm{s}$ end-to-end ($4.4\,\mathrm{GiB}$ peak RSS), which is $10.4\times$ slower than single-threaded GMP \texttt{mpz\_mul} ($0.6056\,\mathrm{s}$).

In single runs on the same machine and operands, the plain-Gaussian variant with the multiplicative schedule (\texttt{algo=0}) took $1.87\times$ and $1.77\times$ as long and used $3.6\times$ and $2.4\times$ the memory at $10^7$ and $10^8$ bits. GMP's \texttt{mpz\_mul} uses a Sch\"onhage--Strassen FFT at these sizes; the constant factors of spreading, twisting, and deconvolution dominate at every tested size, and no crossover is expected at any physically storable size. Accordingly, no practical speedup is claimed.
